% SGA 3，第 VII_B 篇，英文重建稿——第 1 轮。
% 范围：权威底本分篇第 53--54 页，第 2.4.B 节全文，止于命题 2.4.1 之前。
% 合并阅读本页码： 566--567.
% 再版印刷页码： 555--556.
% 权威底本： Polo--Gille Exp7B-13oct24.pdf,
% SHA-256 D303690E0195654BBE7C3DAA5D89EA23E24CD709A3575217234AC19AE8836694.
% 曾参考 Jacob Reinhold 的 CC BY 4.0 Markdown 作为比对谱系。
% 编者注 92 附于 2.4.A，并明确涵盖 2.4.A 与 2.4.B；此处不重复收录。
% 原英文 Loop 1 阶段说明：唯一的图采用直接取自 PDF 的精确源裁片，普通数学均为原生 LaTeX。当前中文工程已采用原生图。

\subsubsection*{2.4.B}
\phantomsection
\label{sga3:VIIB:2.4.B}

为简便起见，假设 \(k\) 为 Artin 环。于是上述讨论可以借助 \(G\)、\(Q\) 和 \(N\) 的协变双代数作更简单的表述。事实上，由于作为形式 \(k\)-簇有 \[
  G=N\times Q
\]，故作为 \(k\)-余代数有 \[
  \mathbf H(G)
  =
  \mathbf H(N)\otimes_k\mathbf H(Q)
\]。此外，由于 \(G\) 中的乘法由 \[
  (x,q)\cdot(x',q')
  =
  \bigl(x\alpha(q,x'),qq'\bigr),
  \qquad
  \alpha(q,x')=qx'q^{-1},
\] 给出，对 \(x\in\mathbf H(N)\) 和 \(q\in\mathbf H(Q)\)，\(\mathbf H(G)\) 中的乘法由 \[
  (x\otimes q)\cdot(x'\otimes q')
  =
  x\phi(q,x')\otimes qq',
\] 给出，其中 \[
  \phi:
  \mathbf H(Q)\otimes_k\mathbf H(N)
  \longrightarrow
  \mathbf H(N)
\] 是由 \(\alpha\) 诱导的 \(k\)-余代数态射。

态射 \(\alpha\) 是下列复合。这里 \(\delta_G\) 和 \(m_G\) 分别表示 \(G\) 的对角态射和乘法态射，\(c_Q\) 为 \(Q\) 中的取逆运算，且 \[
  v(q\otimes q'\otimes x)=q\otimes x\otimes q'.
  % SGA3-VII-DIAG-111
\]。\begin{center}\phantomsection \label{sga3:VIIB:2.4.B:conjugation-composite} % SGA3-VII-DIAG-111
% Authority: Exp7B-13oct24.pdf, local p. 54.
% Crop: figures/Exp7B_localp054_semidirect_conjugation_composite.png
% Locator: tex/components/44_expose_VIIB_section24B_covariant_semidirect_products_en.tex, Loop 1 line 64.
\(
Q\times N
\xrightarrow{\ v\circ(\delta_G\times\operatorname{id})\ }
Q\times N\times Q
\xrightarrow{\ \operatorname{id}\times\operatorname{id}\times c_Q\ }
Q\times N\times Q
\xrightarrow{\ m_G\ }G,
\)
\end{center}也用 \(c_Q\) 表示 \(\mathbf H(Q)\) 的对极，得到 \[
  \phi(q\otimes x')
  =
  \sum_i q_i x'c_Q(q_i')
  \quad\text{if}\quad
  \Delta_{\mathbf H(Q)}(q)
  =
  \sum_i q_i\otimes q_i'.
  \tag{\(\star\)}\label{sga3:VIIB:2.4.B:equation:covariant-antipode-star}
\]。

特别地，设 \(M\) 为抽象群，并假设 \(\mathbf H(Q)\) 为群代数 \(kM\)，从而 \[
  Q=\operatorname{Spf}^{*}kM
\] 为常形式 \(k\)-群 \(M_k\)。则对每个 \(\gamma\in M\) 和 \(x'\in\mathbf H(N)\)，有 \[
  \phi(\gamma\otimes x')
  =
  \gamma x'\gamma^{-1}.
\]。这定义了 \(M\) 通过 Hopf 代数自同构在 \(\mathbf H(N)\) 上的作用。
